\worksheettitle
  {Классы многозначных чисел}
  {\modulemark{M03} Учебный модуль \textbullet\ полный маршрут}

\studentnameblock

\begin{notebox}[Чему вы научитесь]
  Делить запись числа на классы, называть группы классов и количество их
  единиц, а также собирать число по заданным классам.
\end{notebox}

\begin{checkpointbox}[Вход: связь с M02]
  Соберите число и сохраните все разряды:
  \[
    60\,000+3\,000+40+2=\answerblank[30mm].
  \]
  Если получилось не \(63\,042\), сначала выполните коррекционную карточку
  M02-R.
\end{checkpointbox}

\section{Зачем нужны классы}

Длинную запись числа легче читать, если разделить её на группы по три цифры.
Такие группы называют \textbf{классами}.

\begin{problembox}[Разобранный пример]
  Начинаем с правого края числа \(7204061\).
  \begin{enumerate}[label=\textbf{Шаг \arabic*.}]
    \item Последние три цифры образуют класс единиц: \(061\).
    \item Следующие три цифры образуют класс тысяч: \(204\).
    \item Слева остаётся класс миллионов: \(7\).
  \end{enumerate}
  \[
    7204061=7\,|\,204\,|\,061=7\,204\,061.
  \]
\end{problembox}

\begin{reflectionbox}[Правило]
  Запись числа делят на классы \textbf{справа налево}, по три цифры.
  Только крайняя левая группа может содержать одну или две цифры.
\end{reflectionbox}

\newpage
\begin{practicebox}[Пробуем с опорой]
  Отделяйте по три цифры, начиная справа.
  \begin{enumerate}[label=\textbf{\arabic*.}]
    \item \(5308=5\,|\,308=\answerblank[32mm].\)
    \item \(604015=\answerblank[40mm].\)
    \item \(80304072=\answerblank[46mm].\)
  \end{enumerate}
\end{practicebox}

\section{Названия классов}

Классы нумеруют справа налево.

\begin{center}
\renewcommand{\arraystretch}{1.45}
\begin{tabular}{c|ccc}
  \toprule
  номер класса & III & II & I \\
  \midrule
  название & миллионы & тысячи & единицы \\
  \bottomrule
\end{tabular}
\end{center}

\begin{problembox}[Группа класса и количество его единиц]
  Рассмотрим число \(325\,|\,041\,|\,708\).
  \begin{itemize}
    \item Группа класса миллионов: \(325\). Она обозначает \(325\) единиц
      класса миллионов.
    \item Группа класса тысяч: \(041\). Она обозначает \(41\) единицу класса
      тысяч.
    \item Группа класса единиц: \(708\). Она обозначает \(708\) единиц класса
      единиц.
  \end{itemize}
\end{problembox}

\begin{notebox}[Не путайте два разных вопроса]
  В числе \(325\,041\,708\) группа класса тысяч равна \(041\), поэтому в ней
  записана \(41\) единица класса тысяч. Но во всём числе содержится
  \(325\,041\) полная тысяча. В этом модуле мы работаем именно с
  \textbf{группой класса} и \textbf{количеством единиц этого класса}.
\end{notebox}

\begin{notebox}[Почему пишут \(041\), а не \(41\)]
  Каждый класс, кроме крайнего левого, записывают тремя цифрами. Начальные
  нули сохраняют позиции отсутствующих разрядов внутри класса. В группе
  \(041\) ноль показывает отсутствие сотен тысяч.
\end{notebox}

\newpage
\begin{practicebox}[Заполните таблицу]
  Разберите число \(52\,006\,090\).
  \begin{center}
  \renewcommand{\arraystretch}{1.45}
  \begin{tabular}{c|ccc}
    \toprule
    класс & миллионы & тысячи & единицы \\
    \midrule
    группа & \answerblank[15mm] & \answerblank[15mm] & \answerblank[15mm] \\
    количество единиц класса & \answerblank[15mm] & \answerblank[15mm] &
      \answerblank[15mm] \\
    \bottomrule
  \end{tabular}
  \end{center}
\end{practicebox}

\begin{practicebox}[Объясните нули]
  Почему в числе \(52\,006\,090\) нельзя заменить группу \(006\) на \(6\),
  а группу \(090\) на \(90\)?
  \writinglines{2}
\end{practicebox}

\section{Самостоятельная практика: классы числа}

\begin{practicebox}[1. Разделите числа на классы]
  \begin{enumerate}[label=\textbf{\arabic*.}]
    \item \(4072=\answerblank[45mm].\)
    \item \(53008=\answerblank[45mm].\)
    \item \(604015=\answerblank[45mm].\)
    \item \(7030406=\answerblank[52mm].\)
    \item \(80005009=\answerblank[52mm].\)
    \item \(302000070=\answerblank[55mm].\)
  \end{enumerate}
\end{practicebox}

\begin{practicebox}[2. Назовите классы]
  Для числа \(418\,027\,305\) заполните пропуски.
  \begin{enumerate}[label=\textbf{\arabic*.}]
    \item III класс называется классом \answerblank[25mm]; его группа
      \answerblank[18mm].
    \item II класс называется классом \answerblank[25mm]; его группа
      \answerblank[18mm], она обозначает \answerblank[18mm] единиц класса.
    \item I класс называется классом \answerblank[25mm]; его группа
      \answerblank[18mm].
  \end{enumerate}
\end{practicebox}

\begin{practicebox}[3. Найдите нулевой класс]
  В каждом числе подчеркните группу \(000\) и назовите её класс.
  \begin{enumerate}[label=\textbf{\arabic*.}]
    \item \(52\,000\,308\): класс \answerblank[25mm].
    \item \(407\,006\,000\): класс \answerblank[25mm].
    \item \(9\,000\,040\): класс \answerblank[25mm].
  \end{enumerate}
\end{practicebox}

\section{Собираем число по классам}

\begin{problembox}[Разобранный пример: собираем число]
  В числе \(23\) миллиона, \(7\) тысяч и \(5\) единиц.
  \[
    23\;|\;007\;|\;005=23\,007\,005.
  \]
  Группы справа от крайней левой дополняем нулями до трёх цифр.
\end{problembox}

\begin{practicebox}[1. Восстановите числа]
  \begin{enumerate}[label=\textbf{\arabic*.}]
    \item \(36\) миллионов, \(5\) тысяч, \(8\) единиц:
      \(36\;|\;\answerblank[15mm]\;|\;\answerblank[15mm]
      =\answerblank[35mm].\)
    \item \(4\) миллиона, \(0\) тысяч, \(72\) единицы:
      \answerblank[58mm].
    \item \(205\) миллионов, \(9\) тысяч, \(40\) единиц:
      \answerblank[58mm].
  \end{enumerate}
\end{practicebox}

\section{Разбираемся в ошибках}

\begin{warningbox}[2. Исправьте ошибки]
  \begin{enumerate}[label=\textbf{\arabic*.}]
    \item \(7204061=720\,|\,406\,|\,1\).
    \item \(53008=53\,|\,08\).
    \item \(42\) миллиона, \(7\) тысяч, \(90\) единиц записали как
      \(42\,7\,90\).
  \end{enumerate}
  Запишите верные варианты и объясните каждую ошибку.
  \writinglines{4}
\end{warningbox}

\begin{practicebox}[3. Сравните]
  Поставьте знак \(>\), \(<\) или \(=\). Подчеркните первый слева класс, в
  котором группы различаются.
  \[
    8\,041\,005\ \answerblank[8mm]\ 8\,401\,005,
    \qquad
    27\,006\,050\ \answerblank[8mm]\ 27\,060\,005.
  \]
\end{practicebox}

\newpage
\section{Контрольная точка}

\begin{checkpointbox}[Выполните без подсказки]
  \begin{enumerate}[label=\textbf{\arabic*.}]
    \item Разделите на классы: \(90305008=\answerblank[55mm].\)
    \item В полученном числе назовите группы классов миллионов, тысяч и
      единиц.
      \writinglines{1}
    \item Какое количество единиц записано в группе класса тысяч и в группе
      класса единиц?
      \writinglines{1}
    \item Запишите число, в котором \(17\) миллионов, \(4\) тысячи и \(6\)
      единиц: \answerblank[48mm].
    \item Почему в записи этого числа нужны группы \(004\) и \(006\)?
      \writinglines{2}
    \item Исправьте запись \(6150302=615\,|\,030\,|\,2\) и объясните ошибку.
      \writinglines{2}
  \end{enumerate}
\end{checkpointbox}

\begin{reflectionbox}[Проверяю себя]
  \choicebox\ Начинаю деление справа.
  \quad
  \choicebox\ Сохраняю по три цифры во всех группах, кроме крайней левой.

  \choicebox\ Различаю группу класса и количество его единиц.
  \quad
  \choicebox\ Умею собрать число по классам.
\end{reflectionbox}

\begin{notebox}[Дальнейший маршрут]
  Если деление начиналось слева, в группе потеряны нули или перепутаны группа
  класса и количество его единиц, выполните M03-R. Если способ понятен, но
  нужна тренировка, выполните M03-H. Если контроль выполнен без подсказки и
  ошибки объяснены, переходите к M04; M03\(\star\) можно выполнить дополнительно.
\end{notebox}
